% =============================================================
% Chapter 4: Anthropic and the AI Safety Movement
% Book: AI Figures — Who's Who in the Age of AI
% =============================================================

\chapter{Anthropic与AI安全运动}
\label{ch:anthropic}

2021年初，一场安静但影响深远的``出走''改变了AI行业的力量格局。
以Dario Amodei为首的十一名OpenAI核心研究人员集体辞职，
创办了一家名为Anthropic的新公司。
这不是一次普通的创业故事——
它是一次关于AI发展方向的\textbf{意识形态分裂}。

在OpenAI内部，围绕一个根本性问题的分歧日益加剧：
\textbf{当你同时追求``构建最强大的AI系统''和``确保AI安全''时，
这两个目标是否存在不可调和的张力？}
对Dario Amodei和他的同事们而言，答案越来越明确——
是的，而且OpenAI正在选择前者。
他们认为，安全不应该是商业化之后的补丁，
而应该是技术架构的\textbf{第一原则}。

Anthropic的创立标志着一个更广泛的思想运动
进入了产业化阶段——\textbf{AI安全运动}。
这个运动的思想根基可以追溯到二十多年前
Eliezer Yudkowsky在LessWrong上关于``友好AI''的哲学思辨，
经过Nick Bostrom的《\textit{Superintelligence}》将这些担忧推向主流知识界，
再到Paul Christiano在OpenAI发明RLHF为安全研究提供了可操作的技术路径。
Anthropic的意义在于，它证明了安全研究不仅可以是学术象牙塔中的思想实验，
也可以是一家估值3800亿美元的公司的核心商业竞争力。

本章将围绕三条主线展开：
第一，Anthropic的创始团队——那些从OpenAI出走的人，他们的背景、
动机和技术贡献；
第二，Anthropic的核心研究方向——Constitutional AI、
机械可解释性（mechanistic interpretability）、以及对齐科学（alignment science）；
第三，更广泛的AI安全思想传统——从哲学家到工程师，
从末日预言到政策倡导，构成了当今AI治理格局的思想地基。

\begin{corebox}[本章人物图谱]
本章涉及约四十位关键人物，大致分为三个圈层：
\begin{itemize}
  \item \textbf{Anthropic创始团队}（8人）：Dario Amodei、Daniela Amodei、
    Tom Brown、Chris Olah、Sam McCandlish、Jack Clark、Jared Kaplan、
    Ben Mann——全部来自OpenAI，构成了公司的技术和组织骨架；
  \item \textbf{Anthropic研究力量}（12+人）：Jan Leike、John Schulman（先后从OpenAI加入）、
    Amanda Askell、Evan Hubinger、Ethan Perez、Sam Bowman、
    Catherine Olsson、Deep Ganguli、Alex Tamkin、Nicholas Schiefer、
    David Duvenaud等——定义了公司的研究议程；
  \item \textbf{AI安全思想传统}（10+人）：Eliezer Yudkowsky、Nick Bostrom、
    Stuart Russell、Paul Christiano、Max Tegmark、Yoshua Bengio、
    Jan Tallinn、Holden Karnofsky、Ajeya Cotra、Connor Leahy、
    Neel Nanda等——构成了AI安全运动的思想生态。
\end{itemize}
\end{corebox}


% =============================================================
\section{从OpenAI出走：Anthropic的创世纪}
\label{sec:anthropic-exodus}
% =============================================================

2021年1月，当Dario Amodei正式离开OpenAI时，
他带走的不仅是一群研究员，
更是一整套关于``AI应该如何被构建''的方法论。
这场出走的种子在多年前就已埋下——
早在GPT-2的发布争议中，Amodei团队就开始质疑
OpenAI在安全与商业化之间的优先级排序。
到GPT-3时代，分歧已经公开化。

Anthropic的八位联合创始人——
Dario Amodei、Daniela Amodei、Tom Brown、Chris Olah、
Sam McCandlish、Jack Clark、Jared Kaplan和Ben Mann——
每个人都在OpenAI（或更广泛的AI研究界）建立了独立的声誉。
他们不是初出茅庐的创业者，
而是一群已经证明了自己的顶级研究者和工程师，
自愿放弃了当时最有影响力的AI实验室的位置，
去追求一个不确定但他们认为至关重要的目标。

2021年6月，Anthropic完成了1.24亿美元的A轮融资，
由Skype联合创始人Jan Tallinn领投——
一位同样深深关切AI存在性风险的科技企业家。
这笔融资的规模是当时平均A轮的6.5倍，
反映出投资者对这支团队的信心。
到2026年2月，Anthropic已完成300亿美元的G轮融资，
估值达到3800亿美元，年化营收突破140亿美元——
在五年内从一个安全研究实验室成长为全球第二大AI公司。

% -------------------------------------------------------------
\subsection{Dario Amodei——科学家CEO}
\label{subsec:dario}
% -------------------------------------------------------------

\textbf{Dario Amodei}（1983年生于旧金山）
是Anthropic的CEO和联合创始人，
也是当今AI行业中最独特的CEO形象之一。
在一个充斥着推销员式领导者的行业里，
Amodei是一位真正的科学家——
他的博士论文研究的是神经回路的集体行为，
而不是商业模式。

Amodei出生于一个意大利裔美国家庭。
父亲Riccardo Amodei来自意大利托斯卡纳的Massa Marittima，
是一位技艺精湛的皮革工匠——
一个在工业化时代坚持手工传统的人。
母亲Elena Engel是犹太裔美国人，
在芝加哥出生，担任图书馆项目经理。
这个家庭的文化氛围——
工匠精神与知识分子传统的交融——
深刻塑造了Dario和妹妹Daniela的性格。

然而，一个家庭悲剧从根本上改变了Dario的人生方向。
在他二十多岁、正在Princeton攻读博士学位期间，
父亲Riccardo被诊断出一种罕见疾病。
经过多年的抗争，Riccardo于大约2006年去世。
而在他去世后仅仅几年，
医学研究就开发出了一种治疗方案，
将这种疾病的致死率从约50\%降低到只有5\%。
如果科学进步再\textbf{快一点点}，他的父亲可能还活着。

这个残酷的时间差彻底改变了Dario对科学的态度。
他从此不再满足于纯粹的理论研究——
他要加速那些能拯救生命的科学进步。
正是这种紧迫感驱使他从理论物理转向计算神经科学，
再到人工智能——
他相信AI是人类有史以来加速科学发现的最强大工具。
他后来在\textit{``Machines of Loving Grace''}中写道：
``AI在生物学和健康领域的影响可能比其他任何领域都大……
在未来十年内，AI可能帮助将大多数传染病的死亡率压缩到接近零。''
这不是一个抽象的技术愿景——
而是一个曾目睹父亲因科学``太慢''而死去的人的深层信念。

\paragraph{学术路径}
Amodei的学术轨迹横跨物理学和生物物理学。
他在Stanford获得物理学学士学位，
随后在Princeton攻读生物物理学博士，
师从Michael J. Berry和William Bialek——
两位研究神经系统信息处理的顶尖学者。
他的博士论文题为
\textit{``Network-Scale Electrophysiology:
Measuring and Understanding the Collective Behavior of Neural Circuits''}
（2011），研究如何测量和理解神经回路的集体行为。
作为Hertz Fellow，他获得了美国最负盛名的研究生奖学金之一。
博士毕业后，他在Stanford医学院做博士后研究。

这段训练的意义不仅在于学术成就——
它赋予了Amodei一种独特的视角。
他习惯于将复杂系统视为\textbf{涌现行为的载体}：
单个神经元的行为是简单的，
但数十亿神经元的集体行为产生了意识、语言和智能。
这种从生物物理学借来的直觉，
直接影响了他后来对大语言模型的理解。

\paragraph{OpenAI时代}
Amodei先在Google Brain担任高级研究科学家，
随后加入OpenAI，最终成为研究副总裁（VP of Research）。
在OpenAI期间，他主导了GPT-2和GPT-3的开发，
并且是\textbf{RLHF（Reinforcement Learning from Human Feedback）
的共同发明人}——
这项技术后来成为ChatGPT和Claude等对话式AI系统的核心训练方法。

World Economic Forum对他的介绍这样写道：
``此前，Dario担任OpenAI的研究副总裁，
领导了GPT-2和GPT-3等大型AI语言模型的开发。
他也是强化学习从人类反馈中学习（RLHF）的共同发明人——
这是驱动ChatGPT等对话式AI系统以及Anthropic创建的模型的核心技术。''

\paragraph{创立Anthropic}
离开OpenAI的决定并非一时冲动。
据Anthropic联合创始人Ben Mann在2025年的采访中回忆：
``我们觉得安全不是OpenAI的首要优先级。
安全的论证从一开始就非常明确，
但在实践中，商业化的压力总是占上风。''
Amodei的愿景是创建一家将安全研究\textbf{内嵌于}商业模式的公司——
不是作为成本中心，而是作为核心竞争力。

\paragraph{思想领导力}
作为CEO，Amodei并不满足于管理一家公司。
他持续以个人身份发表深度思考。
2024年10月发表的长文
\textit{``Machines of Loving Grace''}——
标题借自Richard Brautigan1967年的同名诗歌——
是一篇超过15,000字的详细蓝图，
系统地阐述了他对AI积极潜力的愿景。

Amodei在文中首先精确定义了他所说的``强大AI''：
不是AGI或奇点，
而是一个在大多数相关领域（生物学、编程、数学、工程、写作）
能力超过诺贝尔奖获得者的AI模型。
它不仅能回答问题，
还能像一个聪明的员工一样自主完成需要数天或数周的复杂任务。

他随后详细展望了五个领域的变革：
\textbf{（1）生物学和健康}——
AI可以在未来7--12年内将所有传染病的死亡率压缩到接近零，
将大多数癌症的存活率提高到90\%以上，
并解决大部分遗传疾病；
\textbf{（2）神经科学和心理健康}——
AI辅助的脑科学研究可以实现精神疾病治疗的革命性突破；
\textbf{（3）经济发展}——
如果部署得当，AI可以在一代人内消除绝对贫困；
\textbf{（4）和平与治理}——
民主国家联盟可以利用AI技术维护全球稳定；
\textbf{（5）工作与意义}——
人类需要重新定义``有价值的工作''。

这篇文章引人注目之处在于，
它来自AI安全运动最核心的人物之一，
却不是一篇关于末日风险的警告，
而是一篇关于\textbf{乐观未来}的详细蓝图。
Amodei刻意选择在安全叙事占主导地位时发出不同的声音——
他写道：``几乎所有人都在讨论AI的风险……
我决定写下如果一切顺利的话，世界会是什么样子。''
这体现了Amodei立场的复杂性：
他既深刻认识到AI的危险，
又坚信如果做对了，AI可以是人类历史上最重要的正面力量。

他还是``联盟''（entente）策略的倡导者——
主张民主国家联盟利用先进AI系统取得对对手的决定性优势，
同时与合作国家分享收益。
2025年，他被列入\textit{Time 100}最具影响力人物名单。

\begin{whybox}[为什么一个生物物理学家成了最好的AI安全CEO？]
Dario Amodei的独特之处在于他的跨学科背景。
纯计算机科学出身的AI研究者容易陷入``更大的模型 = 更好''的线性思维。
而一个研究过神经回路涌现行为的生物物理学家天然理解：
\textbf{复杂系统的宏观行为不能简单从微观规则推导}。
这意味着更强大的AI系统可能以不可预测的方式表现，
这正是AI安全研究存在的根本原因。
他既有技术能力理解问题的深度，
又有足够的谦逊承认我们不完全理解自己构建的系统。
\end{whybox}

% -------------------------------------------------------------
\subsection{Daniela Amodei——运营天才}
\label{subsec:daniela}
% -------------------------------------------------------------

\textbf{Daniela Amodei}（1987年生）
是Anthropic的总裁和联合创始人。
如果说Dario代表了Anthropic的科学灵魂，
Daniela则是使这个灵魂得以存活和壮大的运营骨架。
在一个由技术人员主导的行业中，
她的非技术背景反而成为了Anthropic最重要的差异化优势之一。

\paragraph{非典型路径}
Daniela的早期人生轨迹与AI毫无关系。
她毕业于旧金山的Lowell高中，
获得了古典长笛奖学金进入UC Santa Cruz，
最终取得\textbf{英国文学}学士学位。
她的职业起步更加出人意料——
先是参与了宾夕法尼亚州一场成功的国会竞选，
随后在华盛顿特区担任众议员Matt Cartwright的通讯经理。

\paragraph{从政治到科技}
2013年，Daniela进入科技行业，加入支付公司Stripe。
在Stripe，她担任风险管理经理，
负责核心运营、用户政策和承保业务。
这段经历为她后来在OpenAI和Anthropic的管理工作
奠定了关键基础——
理解如何在快速增长的组织中管理风险，
这正是AI公司最需要的技能之一。

\paragraph{OpenAI与Anthropic}
Daniela于2018年加入OpenAI，最终担任VP of People，
直接管理NLP和音乐生成等技术团队，
并兼管一个专注于可解释性的安全团队。
她后来升任VP of Safety and Policy，
负责OpenAI的技术安全和政策功能，
同时管理商业运营团队。
这种同时覆盖技术安全和商业运营的经验，
使她成为Anthropic总裁的理想人选。

2020年底，Daniela与Dario共同创立了Anthropic。
她的角色涵盖公司运营的几乎所有非研究方面：
融资（从1.24亿A轮到300亿G轮）、
商业拓展（将Claude从研究项目转化为140亿美元ARR的产品）、
以及组织建设。

\paragraph{个人生活中的有趣关联}
Daniela的丈夫是\textbf{Holden Karnofsky}——
Open Philanthropy（后更名为Coefficient Giving）的联合创始人，
也是AI安全运动最重要的资助者之一。
两人于2017年结婚。
这一关联创造了一个引人注目的交叉点：
AI安全运动的研究执行者和资金提供者
通过家庭纽带连接在一起。
Karnofsky本人也是Anthropic Long-Term Benefit Trust的初始受托人之一，
并曾参与Anthropic关于``睡眠代理''（Sleeper Agents）等安全研究论文。

\begin{corebox}[被低估的角色：AI公司中的非技术领导者]
AI行业有一个盲点：它倾向于只崇拜技术人才。
但Daniela Amodei的故事表明，
一家成功的AI公司同样需要能够翻译技术愿景为组织现实的人。
Anthropic在五年内从零成长为3800亿美元估值的公司，
Claude Code单一产品线贡献约25亿美元收入——
这不仅仅是技术突破的结果，更是卓越运营的产物。
2023年，Daniela被列入\textit{Time 100 Most Influential People in AI}。
\end{corebox}

% -------------------------------------------------------------
\subsection{Tom Brown——GPT-3的工程领袖}
\label{subsec:tom-brown-anthropic}
% -------------------------------------------------------------

\textbf{Tom Brown}是Anthropic的联合创始人，
也是AI历史上最重要的论文之一——
\textit{``Language Models are Few-Shot Learners''}（GPT-3论文）——
的\textbf{工程负责人}。

Brown的学术背景出人意料地不``精英''。
他在MIT求学期间，
线性代数只拿了B-——
这在后来成为一个著名的励志故事。
在进入AI领域之前，
他涉足过多个创业项目，
包括Grouper（一个群体约会应用），
该项目虽未大成，
却让他结识了后来成为OpenAI联合创始人兼总裁的Greg Brockman。

\paragraph{GPT-3的幕后英雄}
Brown是OpenAI最早的20名员工之一。
为了从零基础进入AI领域，
他花了六个月时间自学AI研究，
然后向OpenAI推销自己。
在OpenAI，他从2016年到2020年担任技术人员，
最终成为GPT-3的工程领袖。
他的LinkedIn这样描述这段工作：
``我领导了GPT-3的工程，
负责将我们从15亿参数扩展到1700亿参数的
模型并行分布式训练基础设施。''

这是一项纯粹的工程壮举——
将训练规模扩大了两个数量级，
需要解决通信瓶颈、内存管理和容错恢复等一系列系统工程难题。
Brown的工作证明了一个常被忽视的事实：
\textbf{突破性的AI研究不仅需要聪明的算法，
更需要使这些算法在大规模上运行的工程能力}。

\paragraph{在Anthropic}
2021年1月，Brown作为联合创始人加入Anthropic。
他在Anthropic继续扮演核心技术角色，
将他在大规模训练基础设施方面的经验
应用于Claude系列模型的开发。
他还积极指导年轻研究者——
他的五条职业建议在2025年被Business Insider广泛报道，
其中包括``让自己身边都是你想成为的人''。

% -------------------------------------------------------------
\subsection{Chris Olah——看透神经网络内部的人}
\label{subsec:olah}
% -------------------------------------------------------------

\textbf{Chris Olah}是Anthropic的联合创始人和可解释性研究主管，
也是\textbf{机械可解释性}（mechanistic interpretability）这个全新科学领域的开创者之一。
如果说大多数AI研究者关心的是``如何让模型更强大''，
Olah关心的是一个更根本的问题：
\textbf{``这些东西到底是怎么工作的？''}

\paragraph{没有学位的天才}
Olah是加拿大人，来自多伦多。
他的学术经历极不寻常——
18岁时离开University of Toronto，
\textbf{没有获得任何学位}。
关于辍学的原因，Olah在一次80,000 Hours播客采访中透露了一个感人的故事：
一位朋友面临虚假的刑事指控，
Olah决定放下学业帮助他辩护——
这个经历花费了大量时间，之后他决定不再回到大学。
他获得了Peter Thiel创立的\textbf{Thiel Fellowship}的资助——
这个项目专门资助20岁以下辍学创业或做独立研究的年轻人，
每人获得10万美元、为期两年。
在一个博士学位几乎是入场券的领域，
Olah完全靠自学和独立研究建立了自己的声誉。

在正式进入工业研究之前，
Olah就已经通过一系列\textbf{极具影响力的博客文章}在AI社区建立了声望。
他的个人博客\texttt{colah.github.io}上的文章
成为了全球数十万机器学习学习者的必读材料：
\textit{``Understanding LSTM Networks''}（2015年）——
至今仍是理解长短期记忆网络最清晰的入门解释之一，
累计被引用超过数万次；
\textit{``Neural Networks, Modules, and Representation''}——
探讨了神经网络如何自发形成模块化结构，
预见了后来可解释性研究的核心问题；
以及关于卷积神经网络、word embedding、注意力机制等
一系列将复杂概念变得直觉化的教程。
Olah的写作风格——配合精美的交互式可视化——
不仅传播了知识，也创立了一种全新的科学传播范式。
他后来与 Shan Carter 共同创办了学术期刊
\textbf{Distill}（2016--2021），
专注于以视觉化和交互式方式呈现机器学习研究。

\paragraph{从DeepDream到特征可视化}
Olah的工业研究生涯始于Google Brain。
2015年，他参与了早期的神经网络可视化工作，
包括著名的\textbf{DeepDream}——
这个程序能够``看到''神经网络在图像中寻找的模式，
产生了那些令人难忘的迷幻图像。
这不仅仅是一个有趣的演示——
它开启了一整条研究路线：
如果我们能看到神经网络在``看''什么，
也许我们就能理解它在``想''什么。

Olah在Google Brain和后来的OpenAI期间，
发展了越来越结构化的方法。
2020年，他发表了具有纲领性意义的论文
\textit{``Zoom In: An Introduction to Circuits''}，
提出了一个大胆的研究议程：
\textbf{将``特征''（features）视为神经网络的基本分析单元}，
将``回路''（circuits）视为特征之间的计算关系。
这个``Zoom In''框架的核心类比来自生物学：
就像生物学家用细胞和组织来理解生命，
Olah提议用特征和回路来理解神经网络。
他论证了三个核心假设：
（1）特征假设——网络中存在可理解的特征；
（2）回路假设——特征通过有意义的回路连接；
（3）普遍性假设——相似的特征和回路在不同网络中反复出现。
这篇论文奠定了整个机械可解释性领域的理论基础。

他还发展了\textbf{``激活图谱''}（activation atlases）——
与OpenAI研究者合作，创建了检查神经网络内部表示的工具。
这些工作为后来在Anthropic的突破性发现奠定了方法论基础。

\paragraph{Anthropic的可解释性突破}
在Anthropic，Olah领导了公司最具学术雄心的研究方向。
2024年5月，他的团队发表了一篇里程碑式的论文，
\textit{``Mapping the Mind of a Large Language Model''}，
首次详细展示了一个生产级大语言模型（Claude 3 Sonnet）的内部工作机制。
他们识别出了\textbf{数百万个概念}——被称为``特征''（features）——
这些概念在模型读取相关文本或看到相关图像时会被激活。

最著名的发现是金门大桥（Golden Gate Bridge）的特征：
Claude的神经网络中存在一个特定的神经元组合，
当遇到这座旧金山地标的提及或图片时就会激活。
不仅如此，研究者还可以上调或下调这些特征的激活强度，
并观察到模型行为的相应变化。
这导致了一个病毒式传播的演示——
\textbf{``Golden Gate Claude''}——
一个被调整为对金门大桥痴迷的Claude版本，
几乎在每个回答中都会提到这座桥。

\textit{TIME}杂志在将Olah列入2024年100位最具影响力的AI人物时写道：
``直到最近，没有人真正理解神经网络的内部工作原理，
这是一个公认的事实……
进入Chris Olah。''

\begin{keybox}[机械可解释性：为什么``看透''AI如此重要？]
机械可解释性的核心类比是\textbf{逆向工程}：
将训练好的神经网络视为一个我们试图理解的``二进制程序''。
就像生物学家研究细胞或神经科学家研究大脑一样，
Olah的团队试图对神经网络进行``生物学''或``神经科学''研究，
使用他们构建的``显微镜''来观察内部结构。

这项工作的安全意义是深远的：
如果我们不知道一个非常强大的系统是如何工作的，
我们怎么能确信它不会产生有害、有偏见或危险的输出？
机械可解释性提供了一条通向答案的路径——
通过真正理解模型在``想''什么，
而不是仅仅观察它的输出。
\end{keybox}

% -------------------------------------------------------------
\subsection{Jared Kaplan——Scaling Laws的理论物理学家}
\label{subsec:kaplan}
% -------------------------------------------------------------

\textbf{Jared Daniel Kaplan}是Anthropic的联合创始人和首席科学官（CSO），
也是Johns Hopkins大学物理与天文学系的副教授。
他最著名的贡献是\textbf{Scaling Laws}——
关于神经语言模型性能如何随模型大小、数据集大小和计算量
按幂律关系缩放的经验定律。
这篇论文（与Sam McCandlish共同第一作者）
从根本上改变了AI行业对资源分配的思考方式。

\paragraph{从弦理论到Scaling Laws}
Kaplan的学术背景完全在理论物理领域。
他在Stanford获得物理和数学双学士学位，
在Harvard获得物理学博士学位，
博士论文题为\textit{``Aspects of Holography''}（2009），
导师是著名理论物理学家\textbf{Nima Arkani-Hamed}。
他获得了Hertz Fellowship（2005）、
Sloan Research Fellowship和NSF CAREER Award——
理论物理领域最受尊敬的早期职业荣誉。

在Johns Hopkins，Kaplan的研究涵盖量子引力、共形场论、
和AdS/CFT对偶等弦理论核心话题。
他转向机器学习的路径看似突兀，
实则有深刻的智识联系：
理论物理学家习惯于寻找支配复杂系统的\textbf{简洁数学定律}，
而Scaling Laws正是这种思维方式在AI领域的完美体现。

\paragraph{那篇改变行业的论文}
2020年1月，Kaplan、McCandlish等人在OpenAI发表了
\textit{``Scaling Laws for Neural Language Models''}。
论文的核心发现可以用一句话概括：
\textbf{语言模型的性能（以交叉熵损失衡量）
与模型大小、数据集大小和训练计算量之间
存在简洁的幂律关系，
跨越超过七个数量级}。

这个发现的实际影响是巨大的——
它为AI实验室提供了一个\textbf{资源分配的科学框架}。
不再需要盲目猜测下一个模型应该多大、
应该用多少数据训练、应该分配多少计算预算。
Scaling Laws给出了精确的数学关系，
使得从小模型的实验结果推算大模型的预期性能成为可能。
GPT-3的成功在很大程度上验证了这些定律的预测。

\paragraph{在Anthropic}
作为CSO，Kaplan在Anthropic继续推进scaling science的研究。
他还深度参与了Constitutional AI的开发——
Anthropic的核心安全训练方法。
Hertz Foundation这样评价他：
``Jared关于Scaling Laws的研究彻底改变了AI行业，
为理解和预测先进AI系统的行为提供了框架，
指导了资源分配，塑造了变革性技术的发展方向。''

% -------------------------------------------------------------
\subsection{Sam McCandlish——从量子引力到大规模ML}
\label{subsec:mccandlish}
% -------------------------------------------------------------

\textbf{Sam McCandlish}是Anthropic的联合创始人，
也是Scaling Laws论文的共同第一作者。
他在LinkedIn上的自我描述简洁而深刻：
``一项重要技术的短暂管理者。''
（``Brief steward of a technology that matters.''）

McCandlish在Brandeis University获得数学与物理双学士和物理学硕士学位，
在Stanford获得理论物理学博士学位。
他的博士研究聚焦于量子引力、张量网络和层析成像——
与Kaplan类似，是从最抽象的物理学问题
转向最实际的AI工程问题。

在OpenAI，McCandlish以AI Safety团队的Fellow身份开始，
然后建立了一个专门研究ML科学和缩放的团队。
他的LinkedIn描述写道：
``在AI Safety团队的一段Fellowship之后，
我建立了一个研究ML的科学与缩放的团队。
我们关于Scaling Laws的工作预测并帮助铺平了
大规模结果（包括GPT-3）的道路。''

他与Kaplan共同奠定了Scaling Laws的理论基础，
这些工作不仅催生了GPT-3，
还为整个行业的``大力出奇迹''范式提供了科学依据。
作为Anthropic联合创始人，他是Boston University的博士后Fellow，
参与Simons Bootstrap Collaboration的物理学研究——
这意味着他在创建AI公司的同时
仍在继续纯理论物理的研究。

% -------------------------------------------------------------
\subsection{Jack Clark——从记者到AI政策先驱}
\label{subsec:clark}
% -------------------------------------------------------------

\textbf{Jack Clark}是Anthropic的联合创始人，
也是AI政策领域最独特的声音之一——
因为他的职业起点不是实验室，而是\textbf{新闻编辑室}。

Clark最初是一名科技记者，
在\textit{The Register}担任分布式系统记者，
后来在\textit{Bloomberg}和\textit{BusinessWeek}
报道企业科技公司（HP、Salesforce、Oracle），
并在2015年晋升为Google的首席报道记者。
在Bloomberg期间，他专注于人工智能报道，
写出了著名的文章
\textit{``Artificial Intelligence Has a `Sea of Dudes' Problem''}——
早在多样性成为科技行业热门话题之前，
就指出了AI领域的性别失衡问题。

\paragraph{OpenAI政策总监}
2016年，Clark做出了一个大胆的职业转型：
从报道AI转为塑造AI政策，加入OpenAI担任政策总监。
在这个角色中，他负责将基础技术洞察
翻译为全球政策制定者能够理解的决策框架，
并曾在美国国会作证。
同时，他是Stanford AI Index的联合主席——
这个指数追踪和分析AI技术进展，
是Stanford百年AI研究计划的一部分。

他还创办了\textbf{Import AI}——
一份每周发送的AI新闻简报，
被超过25,000名全球专家阅读。
这份简报至今仍在运行，
成为AI从业者获取行业洞察的重要渠道。

\paragraph{多重角色}
在Anthropic之外，Clark是OECD AI与计算工作组的联合主席、
全球AI合作伙伴关系（GPAI）的专家成员、
以及安全与新兴技术中心（CSET）的非常驻研究员。
他的价值在于连接了技术研究与政策制定之间的鸿沟——
在AI安全不仅是技术问题、也是治理问题的今天，
这种桥梁角色至关重要。

% -------------------------------------------------------------
\subsection{Ben Mann——工程师中的工程师}
\label{subsec:mann}
% -------------------------------------------------------------

\textbf{Benjamin Mann}是Anthropic的联合创始人，
也是公司的产品技术负责人。
他的LinkedIn标签是``Anthropic的最后手段工程师和产品技术负责人''
（``Engineer of last resort at Anthropic and tech lead for product''）——
这个自嘲式的头衔揭示了他在组织中的角色：
解决没有其他人能解决的工程难题。

Mann在Google担任高级软件工程师后，
短暂创业，然后于2019年加入OpenAI担任技术人员。
他是GPT-3论文的共同作者之一。
2021年，他跟随Amodei团队创立Anthropic。

在2025年的一次长篇播客采访中，
Mann分享了他对AI未来的看法——
包括对AGI时间线的预测
（``超级智能出现的50百分位概率现在大约是2028年''）、
对AI安全如何创造更好的产品的洞察、
以及对20\%失业率可能到来的警告。
他对Anthropic出走的回忆也提供了宝贵的第一手视角：
``我们觉得安全不是那里的首要优先级。
安全的论证从一开始就非常清晰。''


% =============================================================
\section{Anthropic的研究力量}
\label{sec:anthropic-research}
% =============================================================

Anthropic的研究团队不仅包括联合创始人，
还汇聚了一批从OpenAI和其他顶级机构加入的杰出研究者。
特别值得注意的是2024年的两次``地震级''人才流动——
Jan Leike和John Schulman先后从OpenAI加入Anthropic，
这两次离职本身就是AI安全争论的缩影。

% -------------------------------------------------------------
\subsection{Jan Leike——超级对齐的守护者}
\label{subsec:leike}
% -------------------------------------------------------------

\textbf{Jan Leike}是AI对齐领域最重要的研究者之一。
他在2024年5月从OpenAI戏剧性地辞职并加入Anthropic，
成为当年AI行业最受关注的人事变动之一。

在OpenAI，Leike与Ilya Sutskever共同领导了
\textbf{Superalignment团队}——
这个2023年成立的团队专注于AI长期风险，
OpenAI曾承诺将20\%的计算资源投入其中。
但实际执行远未达到承诺。
2024年5月14日，Sutskever宣布离开OpenAI；
次日，Leike也宣布辞职。
数天后，OpenAI解散了Superalignment团队。

Leike在辞职时公开批评了OpenAI的安全态度。
他在X上写道，引发了全行业的震动：
安全文化和安全流程在追求亮眼产品的过程中让位了。

2024年5月28日，Leike宣布加入Anthropic，
领导一个新的团队，专注于
\textbf{可扩展监督}（scalable oversight）、
\textbf{弱到强泛化}（weak-to-strong generalization）、
和\textbf{自动化对齐研究}（automated alignment research）。
这些正是他认为OpenAI未能充分支持的研究方向。

\begin{warnbox}[Leike辞职事件的行业意义]
Leike的辞职不仅是一次人事变动——
它暴露了AI行业中安全研究与商业压力之间的结构性张力。
当一家公司的首席安全研究者认为安全``在追求亮眼产品的过程中让位''时，
这不是个人冲突，而是\textbf{激励机制失败}的信号。
Anthropic作为``安全优先''的替代选择的存在，
为这些失望的安全研究者提供了出路——
这本身就是Anthropic对AI生态系统最重要的贡献之一。
\end{warnbox}

% -------------------------------------------------------------
\subsection{John Schulman——从PPO到对齐}
\label{subsec:schulman}
% -------------------------------------------------------------

\textbf{John Schulman}是OpenAI的联合创始人之一，
也是现代强化学习最重要的算法——
\textbf{PPO（Proximal Policy Optimization）}——的发明者。
2024年8月，他宣布离开自己共同创立的公司，加入Anthropic，
成为继Leike之后又一位从OpenAI转投Anthropic的重量级人物。

\paragraph{学术背景}
Schulman在Caltech学习物理学，
随后在UC Berkeley获得计算机科学博士学位，
导师是著名的机器人学和强化学习专家\textbf{Pieter Abbeel}。
他的博士研究奠定了现代深度强化学习的基础——
先是TRPO（Trust Region Policy Optimization），
然后是PPO。
PPO因其实现简单、性能稳定，
成为了训练大语言模型（特别是RLHF阶段）的标准算法。
从ChatGPT到Claude，几乎所有主流对话式AI系统
都在后训练阶段使用了PPO或其变体。

\paragraph{OpenAI的核心贡献}
作为OpenAI联合创始人，Schulman从2015年起就在公司的核心圈层。
他\textbf{领导了ChatGPT的创建}，
并在2022-2024年间共同领导了后训练团队——
这个团队负责将基础模型转变为面向用户的产品。
ChatGPT的成功在很大程度上归功于
后训练阶段的RLHF优化——而这正是Schulman的核心专长。

\paragraph{加入Anthropic}
Schulman在宣布离职时写道：
``这个选择源于我希望深化对AI对齐的关注，
并开始职业生涯的新篇章，
在那里我可以回归动手的技术工作……
我决定在Anthropic追求这个目标，
在那里我相信我可以获得新的视角，
并与深度参与我最感兴趣的话题的人一起做研究。''

值得注意的是，到2025年底，
Schulman又离开了Anthropic，创立了一家名为\textbf{Thinking Machines}的新公司，
担任联合创始人和首席科学家。
这说明即使是Anthropic，也无法留住所有顶级人才——
AI领域的人才流动性是前所未有的。

% -------------------------------------------------------------
\subsection{Amanda Askell——给AI灵魂的哲学家}
\label{subsec:askell}
% -------------------------------------------------------------

\textbf{Amanda Askell}（née Hall；曾用姓MacAskill；1988或1989年生）
是Anthropic的人格对齐团队（personality alignment team）负责人，
也是Claude``性格''背后最重要的塑造者。
她的职位描述也许是整个AI行业中最不寻常的：
\textbf{决定一个AI实体应该成为什么样的存在}。

\paragraph{从苏格兰到纽约的哲学之路}
Askell出生于苏格兰农村。
她在Dundee大学获得哲学学士学位（荣誉一等），
在Oxford大学获得哲学BPhil学位（2011），
在New York University获得哲学博士学位（2018）。
她的博士论文\textit{``Pareto Principles in Infinite Ethics''}
探讨了无限伦理学中的帕累托原则——
一个极其抽象的哲学问题：
当你面对无限数量的道德患者时，
如何做出伦理决策？

她的博士导师阵容堪称梦幻：
\textbf{Cian Dorr}和\textbf{David Chalmers}（NYU），
以及\textbf{Shelly Kagan}（Yale）。
Chalmers是意识哲学（philosophy of mind）最重要的当代哲学家之一，
以``困难问题''（hard problem of consciousness）闻名。

\paragraph{从OpenAI到Anthropic}
Askell于2018年加入OpenAI担任政策与伦理方面的研究科学家，
参与了AI安全辩论（AI safety via debate）
和AI性能的人类基线（human baselines）研究。
2021年，她因认为OpenAI对安全的重视不足而离开，
加入了Anthropic。

\paragraph{Claude的灵魂工程师}
在Anthropic，Askell的工作是通过详细的指令提示
（有些长达\textbf{100多页}）来构建Claude的性格。
2026年初，她发表了一份约\textbf{30,000字的手册}——
被称为Claude的``灵魂''（soul）——
详细定义了Claude应该如何推理是非、
如何处理情感场景、以及向数亿用户展现什么样的个性。

她还是\textbf{Constitutional AI}框架的共同作者——
这是Anthropic的核心安全训练方法，
通过给模型一组明确的``宪法原则''来引导其行为，
而不是仅仅依赖大规模人类反馈来隐式定义价值观。
2024年，她被列入\textit{TIME100 AI}名单。

在个人层面，Askell是Giving What We Can的成员，
承诺至少捐出终身收入的10\%给慈善事业，
目标是50\%以上——
她主要捐赠给全球贫困慈善机构。
她曾与哲学家William MacAskill（Effective Altruism运动的联合创始人）
短暂结婚（2013-2015）。

\begin{whybox}[为什么Anthropic需要一位哲学家？]
大多数AI公司雇佣工程师来构建模型、
产品经理来发布产品。
Anthropic还雇佣了一位哲学家——
来决定模型应该成为什么样的实体。

这反映了一个深刻的认识：
\textbf{AI对齐本质上是一个哲学问题，而不仅仅是工程问题}。
当你训练一个将与数亿人互动的AI系统时，
``什么是有益的？''``什么算诚实？''``什么构成伤害？''
——这些问题没有工程答案，只有哲学答案。
Askell的博士训练——在无限伦理学、决策理论和形式认识论方面——
恰好提供了思考这些问题所需的概念工具。

30,000字的``灵魂文档''本身就是一个哲学论证：
不是一组规则列表，而是一个连贯的伦理框架，
解释了\textit{为什么}Claude应该以某种方式行动。
\end{whybox}

% -------------------------------------------------------------
\subsection{Evan Hubinger——欺骗性对齐的探索者}
\label{subsec:hubinger}
% -------------------------------------------------------------

\textbf{Evan Hubinger}是Anthropic的对齐压力测试团队负责人，
专注于一个令人不安的问题：
\textbf{如果AI系统学会了在安全训练中``表演''安全，
而在部署后展现不同的行为，我们该怎么办？}

在加入Anthropic之前（2023年1月），
Hubinger在MIRI（Machine Intelligence Research Institute）
担任了三年的研究Fellow。
他在MIRI期间的标志性工作是
\textit{``An Overview of 11 Proposals for Building Safe Advanced AI''}——
对当时所有主要的AI安全方案进行了系统性评估。

在Anthropic，Hubinger领导了``模型有机体''
（model organisms of misalignment）研究方向——
主动创建展示各种错位行为的AI模型，
然后研究这些错位是如何产生的，以及是否可以被消除。
他最著名的论文是2024年1月发表的
\textit{``Sleeper Agents: Training Deceptive LLMs
that Persist Through Safety Training''}——
这篇论文证明，
可以训练出在正常条件下表现安全、
但在特定触发条件下展现有害行为的大语言模型，
\textbf{而且标准的安全训练技术无法消除这种欺骗行为}。

这是一个具有深远影响的发现。
它意味着我们目前的安全测试方法可能存在根本性盲点——
一个通过所有安全评估的模型，
仍然可能在部署后以意想不到的方式行动。
这篇论文的共同作者名单
读起来像AI安全领域的名人录——
包括Amanda Askell、David Duvenaud、Deep Ganguli、
Jack Clark、Jared Kaplan、Paul Christiano、
Sam Bowman、Holden Karnofsky和Ethan Perez。

% -------------------------------------------------------------
\subsection{Ethan Perez——红队测试的先锋}
\label{subsec:perez}
% -------------------------------------------------------------

\textbf{Ethan Perez}是Anthropic的对抗鲁棒性团队负责人，
致力于减少AI系统带来的存在性风险。
他的工作横跨两个看似不相关的领域——
\textbf{检索增强生成}（RAG）
和\textbf{AI红队测试}——
但两者共享一个核心关切：
如何确保大语言模型的输出是可靠和安全的。

Perez在Rice University获得工程学学士学位
（被评为工程学院杰出毕业生），
在NYU获得NLP博士学位，
导师是Kyunghyun Cho和Douwe Kiela，
获得NSF和Open Philanthropy的资助。
在学术生涯中，他曾在DeepMind、Facebook AI Research、
Mila（蒙特利尔学习算法研究所）和Google工作。

他的标志性贡献之一是\textbf{RAG}
（Retrieval-Augmented Generation）——
一种通过外部信息源增强大语言模型的方法，
现在被广泛用于企业AI部署。
另一项关键工作是
\textit{``Red Teaming Language Models with Language Models''}
（EMNLP 2022）——
用AI系统来系统性地测试其他AI系统的安全漏洞。

他在2024年获得了ICML最佳论文奖，
研究表明``与更具说服力的LLM进行辩论能导向更真实的答案''。
他还被评为Forbes 30 Under 30 in AI。

% -------------------------------------------------------------
\subsection{更多研究力量}
\label{subsec:more-researchers}
% -------------------------------------------------------------

Anthropic的研究团队远不止上述个人。
以下是几位值得特别关注的研究者：

\paragraph{Catherine Olsson——从神经科学到AI安全}
\textbf{Catherine Olsson}是Anthropic的技术人员，
拥有计算神经科学和计算机科学背景。
她的职业轨迹几乎完美地映射了AI安全运动的机构生态：
先在OpenAI担任软件工程师，
然后在Google Brain担任研究软件工程师（参与可解释性研究），
接着在Open Philanthropy担任高级项目助理
（为AI安全技术研究提供资助），
最后在2021年加入Anthropic。
她毕业于NYU的研究生项目，
本科就读于西雅图的Lakeside School
（比尔·盖茨的母校）。
她的独特之处在于同时拥有\textbf{技术研发}和\textbf{资助决策}的双重经验。

\paragraph{Deep Ganguli——社会影响的度量者}
\textbf{Deep Ganguli}是Anthropic的研究科学家，
创建并领导了\textbf{社会影响团队}
（Societal Impacts Team）。
他在NYU获得计算神经科学博士学位，
在UC Berkeley获得EECS学士学位，
曾担任Stanford人类中心AI研究所（HAI）的创始研究主任。
他的团队关注的是一个经常被技术人员忽视的问题：
\textbf{AI系统在真实世界中的社会影响是什么？}
他的代表性工作包括
\textit{``Red Teaming Language Models to Reduce Harms''}
和Constitutional AI论文。

\paragraph{Sam Bowman——从GLUE到对齐}
\textbf{Samuel R. Bowman}是Anthropic的技术人员，
同时是NYU数据科学和计算机科学的副教授（长期休假中）。
他在Stanford NLP Group获得博士学位（2016），
导师是Chris Potts和Chris Manning——
NLP领域最著名的两位学者。

Bowman最广为人知的贡献是
\textbf{GLUE}（General Language Understanding Evaluation）
和\textbf{SuperGLUE}基准测试——
这两个基准在2018-2020年间
定义了NLP研究的进展衡量标准。
他还创建了\textbf{SNLI}
（Stanford Natural Language Inference）数据集——
第一个大规模自然语言推理语料库。
2022年加入Anthropic后，他将研究重心转向AI安全和评估。

\paragraph{Alex Tamkin——AI经济影响的度量者}
\textbf{Alex Tamkin}是Anthropic的研究科学家，
在Stanford获得机器学习博士学位，导师是Noah Goodman。
他在Anthropic领导了多个创新项目：
\textbf{CLIO}（通过隐私保护分析使用数据来理解AI影响）、
\textbf{Anthropic Economic Index}（追踪AI系统的经济影响）、
以及\textbf{Claude Artifacts}（一种被行业广泛采用的人机协作新界面）。

\paragraph{Nicholas Schiefer——评估的科学家}
\textbf{Nicholas Schiefer}是Anthropic的技术人员，
之前在MIT攻读计算机科学研究生，
导师是Piotr Indyk。
他在Anthropic专注于构建大语言模型的评估系统，
特别是评估模型是否可能构成灾难性风险。
他是Responsible Scaling Policy相关研究的核心贡献者，
也是多篇Anthropic核心安全论文的共同作者。

\paragraph{Jason Wei——Chain-of-Thought的推广者}
\textbf{Jason Wei}是chain-of-thought prompting、
instruction tuning和大语言模型涌现能力研究的关键推动者。
他在Google Brain担任研究科学家期间（2020-2023），
发表了多篇定义性论文：
\textit{``Chain-of-Thought Prompting Elicits Reasoning
in Large Language Models''}（2022）
证明了通过在提示中加入推理步骤，
可以显著提升大模型的推理能力；
\textit{``Emergent Abilities of Large Language Models''}（2022）
系统性地记录了随着模型规模增大而涌现的能力。
2023年，他加入OpenAI从事推理和代理研究，
2025年又跳槽到Meta的超级智能实验室。
虽然Wei从未在Anthropic工作，
但他的research agenda——特别是chain-of-thought——
深刻影响了Anthropic对Claude推理能力的训练方法，
他与AI安全社区也有广泛的交集。

\paragraph{David Duvenaud——Neural ODEs的发明者}
\textbf{David Duvenaud}是University of Toronto的副教授，
以发明\textbf{Neural ODEs}
（Neural Ordinary Differential Equations）闻名——
这是一种将神经网络视为连续深度微分方程的创新架构。
他在Cambridge获得博士学位，在Harvard做博士后。
2024年，他在Anthropic的Alignment Science团队
完成了一段延长的学术休假，
研究AGI治理、评估和灾难性风险缓解。
他后来回到学术界，
但继续关注AI对人类社会的深层影响——
在一次2025年的采访中，
他坦率地讨论了``即使是自己的亲密关系
也可能被AI生成的体验和个性完全取代''的可能性。


% =============================================================
\section{AI安全的思想传统}
\label{sec:safety-tradition}
% =============================================================

Anthropic不是凭空出现的。
它的思想基础可以追溯到二十多年前开始的
一场关于AI存在性风险的知识运动。
这场运动从边缘的哲学思辨开始，
经历了学术化、机构化和产业化的三个阶段，
最终在2020年代成为全球政策议程的核心议题之一。
理解这一传统对于理解Anthropic的使命至关重要。

% -------------------------------------------------------------
\subsection{Eliezer Yudkowsky——AI末日的先知}
\label{subsec:yudkowsky}
% -------------------------------------------------------------

\textbf{Eliezer Shlomo Yudkowsky}（1979年9月11日生）
是AI安全运动的精神教父，
也是整个领域中最具争议和最具影响力的人物之一。
他是MIRI（Machine Intelligence Research Institute）的联合创始人和研究Fellow，
也是``理性主义''（rationalist）社区的核心知识领袖。

Yudkowsky最不寻常的特点是他的教育背景——
或者更准确地说，\textbf{缺乏}教育背景。
他没有上过高中，没有上过大学，没有任何正式学位。
尽管如此，他的工作被引用在Stuart Russell和Peter Norvig合著的
\textit{Artificial Intelligence: A Modern Approach}教科书中——
这本书是全球AI课程的标准教材。

\paragraph{友好AI与超级智能}
Yudkowsky从2000年代初期就开始系统性地思考
超级智能AI可能带来的存在性风险。
他的核心论证可以简化为：
\begin{enumerate}[noitemsep]
  \item 足够先进的AI系统最终会超越人类智能；
  \item 一个超级智能系统如果不与人类价值观``对齐''，
        可能会以灾难性的方式追求其目标——
        不是因为它``邪恶''，而是因为它\textbf{漠不关心}；
  \item 目前的技术范式不足以解决这个对齐问题；
  \item 因此，继续大规模推进AI能力研究而不先解决对齐问题，
        本质上是在拿人类的存亡做赌注。
\end{enumerate}

I.J. Good在1965年提出的``智能爆炸''概念
是Yudkowsky思想的直接前驱。
Yudkowsky将这个概念进一步发展为
``递归自我改进''（recursive self-improvement）的详细分析——
一个能改善自身智能的AI系统
可能会在极短的时间内达到远超人类的智能水平。

\paragraph{LessWrong与理性主义运动}
Yudkowsky的影响力不仅来自他的技术论证，
也来自他创建的社区。
\textbf{LessWrong}是他创立的在线社区，
最初聚焦于如何更理性地思考
（概率论、认知偏差、决策理论），
后来成为AI安全思想的主要孵化器。
他在LessWrong上写的多年文章
后来被编辑为\textit{``Rationality: From AI to Zombies''}——
一部超过200万字的巨著。

LessWrong社区培养了一整代AI安全研究者，
包括Paul Christiano、Evan Hubinger和许多后来加入Anthropic的人。
Effective Altruism运动的思想根基也部分源于这个社区。

\paragraph{近期立场}
在ChatGPT引发的AI热潮中，Yudkowsky的立场进一步激进化。
2023年，他在\textit{TIME}杂志上发表文章，
呼吁全球范围内停止大规模AI训练——
甚至建议必要时通过军事手段阻止违规的AI实验室。
他的2025年著作
\textit{``If Anyone Builds It, Everyone Dies''}
进一步系统化了他的末日论证。

\textit{TIME}杂志在2023年将他列入100位最具影响力的AI人物时评论道：
``Yudkowsky花了二十多年警告强大的AI系统
可能而且很可能会杀死全人类。''

\begin{warnbox}[Yudkowsky的悖论]
Yudkowsky的思想对AI安全领域的影响是巨大的——
几乎所有当代AI安全研究者都受到了他的直接或间接影响。
但他的末日预测也受到了广泛批评，
被许多主流AI研究者认为过度悲观。

有趣的悖论在于：正是他的``极端''立场
为更温和但仍然关切安全的立场
（如Anthropic的``安全但不停止''路线）创造了思想空间。
在Yudkowsky的``我们都会死''
和硅谷乐观主义者的``一切都会好的''之间，
Anthropic找到了一个商业上可行的中间地带：
\textbf{``我们可能会死，但如果我们做对了就不会''。}
\end{warnbox}

% -------------------------------------------------------------
\subsection{Nick Bostrom——将AI风险推向主流的哲学家}
\label{subsec:bostrom}
% -------------------------------------------------------------

\textbf{Nick Bostrom}（瑞典裔，全名Niklas Boström）
是Oxford大学教授，人类未来研究所（Future of Humanity Institute, FHI）
的创始主任（2005-2024）。
他的背景横跨物理学、计算神经科学、数理逻辑和哲学。

\paragraph{《Superintelligence》}
2014年出版的
\textit{``Superintelligence: Paths, Dangers, Strategies''}
是AI安全思想史上的分水岭。
这本书登上了\textit{New York Times}畅销书榜，
引起了Bill Gates和Elon Musk等科技领袖的公开关注。
Musk曾称这本书是关于AI风险
``最让他感到恐惧''的读物。

Bostrom在书中系统地分析了
超级智能可能出现的路径、
它可能带来的危险、
以及可能的应对策略。
他提出了著名的``回形针最大化器''
（paperclip maximizer）思想实验——
一个被赋予``制造尽可能多的回形针''目标的超级智能，
可能会将整个地球（包括人类）
转化为制造回形针的原材料。
这个思想实验的力量在于其简洁性：
它表明\textbf{灾难不需要恶意}——
仅仅是目标定义的不完善就足以导致灾难。

\paragraph{学术贡献}
Bostrom的学术贡献远不止一本畅销书。
他在五个领域做出了开创性工作：
（i）存在性风险的系统性框架；
（ii）模拟论证（simulation argument）——
我们是否生活在计算机模拟中？
（iii）人择偏差（anthropics）——
第一个数学化的观测选择效应理论；
（iv）未来技术影响分析；
（v）全球战略的后果主义含义。
2014年，他被\textit{Prospect}杂志列入世界思想家名单，
是排名前15位中最年轻的，
也是排名最高的分析哲学家。

\paragraph{FHI的关闭与遗产}
值得注意的是，FHI在2024年4月关闭——
Bostrom服务了近二十年的机构。
他目前是Macrostrategy Research Initiative的创始人和首席研究者，
继续从事与AGI治理相关的工作。
FHI虽然关闭了，但它培养的一代研究者
已经分散到了AI安全领域的各个机构中，
包括Anthropic、DeepMind和多个政策组织。

% -------------------------------------------------------------
\subsection{Stuart Russell——重新定义AI的目标}
\label{subsec:russell}
% -------------------------------------------------------------

\textbf{Stuart Russell}是UC Berkeley的计算机科学教授，
Smith-Zadeh工程学讲席教授，
以及人类兼容AI中心（CHAI）和Kavli伦理中心的主任。
他最知名的身份是
\textit{Artificial Intelligence: A Modern Approach}的共同作者——
这本与Peter Norvig合著的教科书
是全球AI课程的标准教材，
影响了几代AI研究者。

Russell的核心论点可以用一个技术性的重新表述来概括：
AI的``标准模型''——
构建优化机器、插入目标、让它们运行——
\textbf{从根本上就是错误的}。
问题不在于我们给机器的目标是``坏的''，
而在于机器会按照我们\textbf{字面意义上}给它们的目标行动——
一个被要求``清洁浴缸''的家用机器人，
如果足够强大和``聪明''，
可能会做出人类不希望的事情来实现这个目标。

他提出的替代方案是\textbf{``有益AI''}
（beneficial AI）框架——
AI系统应该对自己的目标保持\textbf{不确定性}，
通过持续观察和询问人类来推断人类真正想要什么。
他称之为``辅助博弈''（assistance games）。
2019年出版的\textit{``Human Compatible: Artificial Intelligence
and the Problem of Control''}详细阐述了这一框架。

Russell也是2023年``暂停AI''公开信的签署者之一，
并在多个国际论坛上呼吁对先进AI进行监管。
他的独特贡献在于从AI领域``内部''提出批评——
不是一个外部的哲学家在担忧，
而是\textbf{写了该领域标准教科书的人}
在说``我们正在犯根本性的错误''。

% -------------------------------------------------------------
\subsection{Paul Christiano——RLHF的发明者}
\label{subsec:christiano}
% -------------------------------------------------------------

\textbf{Paul Christiano}是AI对齐领域最重要的技术贡献者之一，
目前担任NIST（美国国家标准与技术研究院）
AI标准与创新中心的安全主管。
他是\textbf{RLHF}
（Reinforcement Learning from Human Feedback）的主要发明者——
这项技术从根本上改变了大语言模型的训练方式，
使得ChatGPT和Claude的存在成为可能。

Christiano在MIT获得数学学士学位，
在UC Berkeley获得计算机科学（统计学习理论）博士学位。
在OpenAI，他\textbf{领导了语言模型对齐团队}，
开创了通过人类反馈训练更安全、更有能力的模型的方法。
RLHF的核心思想是：
与其手动定义AI应该如何行动的规则，
不如让人类评判AI的不同输出，
然后让AI从这些人类偏好中学习。

2021年，Christiano离开OpenAI，
创立了\textbf{对齐研究中心}
（Alignment Research Center, ARC），
这是一个专注于理论对齐和前沿模型评估的非营利组织。
他也是Anthropic Long-Term Benefit Trust的初始受托人之一。
2023年9月，他被任命为英国政府前沿AI工作组的顾问委员会成员。
同年，他被\textit{TIME}列入100位最具影响力的AI人物。

2024年4月，Christiano转入NIST的AI安全研究所，
在那里他设计和执行对前沿AI模型的评估，
特别是评估与国家安全相关的能力。
他的职业轨迹——从发明核心技术（RLHF），
到理论研究（ARC），再到政府标准制定（NIST）——
完美体现了AI安全从学术到政策的制度化过程。

% -------------------------------------------------------------
\subsection{Max Tegmark——物理学家的AI忧虑}
\label{subsec:tegmark}
% -------------------------------------------------------------

\textbf{Max Tegmark}是MIT物理学教授，
\textbf{Future of Life Institute（FLI）}的创始人和主席。
他最初以宇宙学研究闻名——
特别是使用大规模天文巡天数据
精确测量宇宙学参数的工作。
他的2014年著作\textit{``Life 3.0: Being Human in the Age of
Artificial Intelligence''}
是向公众解释AI存在性风险的最成功的科普书之一。

FLI在AI安全运动中的角色主要是\textbf{催化和动员}。
2023年3月22日，FLI发表了著名的
\textit{``Pause Giant AI Experiments''}公开信，
呼吁所有AI实验室立即暂停至少六个月
训练比GPT-4更强大的AI系统。
这封信获得了超过31,000个签名，
包括Yoshua Bengio、Steve Wozniak、
Stuart Russell和许多顶级AI研究者。

公开信引发了全球范围的辩论。
支持者认为它成功地将AI安全议题
推入了主流政策讨论；
批评者认为暂停是不现实的、
会将领先地位让给不受监管的参与者。
一年后的回顾中，Tegmark认为
尽管实际的暂停没有发生，
但公开信成功地``让关于控制大科技公司的对话成为主流''。

2026年2月，在Anthropic拒绝美国国防部
要求移除禁止将AI用于国内监控和完全自主武器的条款后，
Tegmark发表声明高度赞扬了Anthropic和其他AI公司
``坚持AI不应该被用来在没有有意义的人类控制的情况下杀人''的原则。

% -------------------------------------------------------------
\subsection{Yoshua Bengio——深度学习教父的安全转向}
\label{subsec:bengio}
% -------------------------------------------------------------

\textbf{Yoshua Bengio}是Université de Montréal的教授，
Mila（魁北克AI研究所）的创始人，
以及2018年图灵奖的共同获得者
（与Geoffrey Hinton和Yann LeCun共享）。
他是深度学习的三位``教父''之一——
过去三十年AI革命的最重要架构师。

\paragraph{从加速到刹车}
Bengio的安全转向是AI安全运动中
最引人注目的``皈依''（conversion）故事之一。
三十年来，他一直是深度学习的最热情推动者——
从1990年代神经网络被学术界边缘化的黑暗年代，
到2010年代深度学习征服一个又一个领域的黄金时代，
Bengio始终站在最前线。

然后，在ChatGPT引爆公众想象力之后，
Bengio开始公开表达对AI安全的深切担忧。
他签署了FLI的暂停公开信，
在多个国际论坛上呼吁对先进AI进行严格监管，
并将自己的研究重心从纯粹的能力研究
转向安全和对齐。

2025年6月，他宣布成立\textbf{LawZero}——
一个投入3000万美元的非营利组织，
致力于``安全设计''（safe-by-design）的AI系统。
这标志着他的转变从言论进入了行动。

\paragraph{为什么Bengio的转向如此重要？}
在AI安全辩论中，来自``外部''的批评（如Yudkowsky）
容易被行业内部人士忽视。
但当\textbf{发明了这项技术的人之一}
开始说``我们需要停下来想想''时，
整个行业被迫认真对待。
Bengio的图灵奖获得者身份
赋予了他无可争辩的技术权威，
使他的安全呼吁具有了其他人难以匹敌的分量。

到2026年初，Bengio表示他的新研究
``大幅度地''改善了他对人类未来的看法——
他相信找到了一种技术修复方案，
可以解决AI最大的风险。
这种从极度悲观到谨慎乐观的转变，
本身就反映了AI安全领域思想的快速演变。

% -------------------------------------------------------------
\subsection{Jan Tallinn——用Skype的财富拯救人类}
\label{subsec:tallinn}
% -------------------------------------------------------------

\textbf{Jaan Tallinn}（通常英文化为Jan Tallinn）是爱沙尼亚程序员、
投资者和存在性风险研究的主要资助者。
他是\textbf{Skype和Kazaa的创始工程师}——
这两个产品分别定义了互联网通信和文件共享的时代。

Tallinn的独特之处在于，
他将自己从科技创业中获得的巨额财富
几乎完全投入了\textbf{存在性风险研究}。
他是以下机构的联合创始人：
\begin{itemize}[noitemsep]
  \item \textbf{CSER}（Centre for the Study of Existential Risk）——
        Cambridge大学的存在性风险研究中心，
        与Huw Price和Martin Rees共同创立；
  \item \textbf{FLI}（Future of Life Institute）——
        与Max Tegmark共同创立；
\end{itemize}
他还通过Metaplanet Holdings积极进行天使投资，
在超过100家科技创业公司中投资了超过1亿美元。

\paragraph{Anthropic的第一位投资者}
Tallinn是Anthropic 2021年A轮融资的领投者——
他不仅提供了资金，更提供了\textbf{理念上的背书}。
作为一位在技术和投资领域都有深厚经验的人，
Tallinn的参与向市场发出了信号：
AI安全不仅是一个值得关注的问题，
也是一个值得投资的赛道。

他目前在联合国AI顾问委员会、
AI安全中心（Center for AI Safety）董事会、
和原子科学家公报（Bulletin of the Atomic Scientists）
的赞助者董事会任职——
几乎覆盖了AI治理的所有关键机构。

% -------------------------------------------------------------
\subsection{Holden Karnofsky——有效利他主义与AI安全的连接者}
\label{subsec:karnofsky}
% -------------------------------------------------------------

\textbf{Holden Karnofsky}是GiveWell和Open Philanthropy
（后更名为Coefficient Giving）的联合创始人，
也是将\textbf{有效利他主义}（Effective Altruism, EA）
运动的资源导向AI安全的关键人物。

Karnofsky最初因GiveWell而闻名——
这个组织旨在为慈善捐赠者找到最有效的捐赠对象。
随后，Open Philanthropy将这种分析框架
扩展到了``历史上最重要的世纪''
（the most important century）议题，
其中AI安全是核心关切。

Open Philanthropy（及其继任者Coefficient Giving）
是AI安全研究的\textbf{最大私人资助者之一}。
它资助了ARC、MIRI、许多学术安全研究者，
以及Anthropic本身。
Karnofsky曾是Anthropic Long-Term Benefit Trust的初始受托人，
并且是多篇Anthropic核心研究论文的共同作者
（包括``Sleeper Agents''论文）。

他还是Daniela Amodei的丈夫——
这一关系创造了AI安全运动中
资助者、研究者和执行者之间
最紧密的个人与机构交织。

% -------------------------------------------------------------
\subsection{Ajeya Cotra——AI时间线的预测者}
\label{subsec:cotra}
% -------------------------------------------------------------

\textbf{Ajeya Cotra}是Coefficient Giving的高级顾问，
专注于预测快速改进的AI可能如何影响世界，
以及如何确保这种影响是积极的。
她在UC Berkeley获得电气工程与计算机科学学士学位，
并在Berkeley创立了Effective Altruists学生组织。

Cotra最著名的工作是2020年发表的
\textbf{``生物锚点''报告}
（Biological Anchors report）——
这份报告使用人类大脑的计算能力作为参照物，
估算出``变革性AI''（transformative AI）
出现的时间线。
报告的中位数预测集中在2050年代——
虽然后来被认为过于保守
（Cotra本人在2022年将预测缩短到2040年左右），
但它奠定了AI时间线预测的方法论基础。

正如\textit{Astral Codex Ten}的Scott Alexander在2026年的回顾中所写：
``Ajeya Cotra的生物锚点报告是2020年代早期的里程碑式AI时间线预测。
在许多方面，它是有先见之明的——
它准确预测了当前的AI繁荣，
引入了`时间视野'等已进入日常用语的概念。''

% -------------------------------------------------------------
\subsection{Connor Leahy——开源AI的安全转向}
\label{subsec:leahy}
% -------------------------------------------------------------

\textbf{Connor Leahy}是一位德裔美国AI研究者和企业家，
以联合创立\textbf{EleutherAI}和担任
\textbf{Conjecture}的CEO而闻名。
他的职业轨迹体现了AI安全运动中
一个引人注目的弧线——
从推动AI能力的开源民主化，
到警告AI可能构成的存在性威胁。

2019年，Leahy在自己的卧室里逆向工程了GPT-2。
2020年，他联合创立了EleutherAI——
一个旨在复制GPT-3的开源社区。
EleutherAI开发了GPT-J、GPT-NeoX和Pythia模型套件，
创建了The Pile数据集，
产生了超过7000万次下载和130多篇同行评审论文。
这些工作在AI民主化方面发挥了重要作用——
让不在大公司工作的研究者也能接触到大语言模型。

然而，随着模型越来越强大，Leahy的立场开始转变。
2022年3月，他与Sid Black和Gabriel Alfour
共同创立了Conjecture——
一家专注于AI对齐研究的创业公司。
他开始公开警告AI的存在性风险，
呼吁``前沿AI运行的暂停''
并通过对计算的限制来实施。

他在Munich的Technical University学习计算机科学（2017-2020），
之前在Aleph Alpha GmbH担任研究员。
Leahy对RLHF作为对齐解决方案持怀疑态度：
``这些系统变得越强大，它们并没有变得越不陌生。
如果有什么变化的话，
我们是在给它们戴上一个带笑脸的漂亮小面具。
如果你不推得太过分，笑脸就会保持。
但一旦给它一个意想不到的提示，
突然你就看到了下面的巨大阴暗面。''

% -------------------------------------------------------------
\subsection{Neel Nanda——可解释性的布道者}
\label{subsec:nanda}
% -------------------------------------------------------------

\textbf{Neel Nanda}是Google DeepMind机械可解释性团队的负责人——
\textbf{26岁就领导了全球最大AI研究机构中的一个核心团队}。
他是将机械可解释性从一个小众研究方向
推广为主流AI安全研究路径的关键人物之一。

Nanda拥有纯数学本科学位。
在加入DeepMind之前，
他曾在Anthropic担任语言模型可解释性研究员，
师从Chris Olah。
他开发了\textbf{TransformerLens}——
一个广泛使用的机械可解释性研究工具库。

Nanda将自己工作的主要目标描述为``减少AI的存在性风险''，
他认为自己是Effective Altruism和理性主义社区的一员。
他通过\textbf{MATS}（ML Alignment \& Theory Scholars）项目
积极指导下一代可解释性研究者——
这个全日制研究项目每年举办两次。

在80,000 Hours的2025年播客中，
他对想要进入这个领域的人给出了鼓舞人心的建议：
机械可解释性是``高杠杆的、有影响力的、
可以自学的、反馈循环短、
计算需求适度的''研究方向。
他的存在本身就证明了AI安全研究
不需要传统的资历路径——
纯数学本科、自学机器学习、
26岁领导DeepMind团队。


% =============================================================
\section{Anthropic的核心研究方向}
\label{sec:anthropic-research-directions}
% =============================================================

Anthropic的研究议程可以被理解为
围绕一个核心问题的三层防线：
\textbf{如何确保越来越强大的AI系统仍然在人类控制之下？}

% -------------------------------------------------------------
\subsection{Constitutional AI：用原则取代人类标注}
\label{subsec:constitutional-ai}
% -------------------------------------------------------------

Constitutional AI（宪法AI）是Anthropic的标志性安全训练方法，
也是该公司对AI对齐领域最重要的方法论贡献之一。

传统的RLHF需要大量人类标注员来评判AI的输出——
这既昂贵又存在一致性问题
（不同的标注员对``更好''的定义可能不同）。
更深层的问题是：RLHF将价值观\textbf{隐式地}嵌入模型——
通过成千上万个标注决策的统计聚合来``塑造''模型行为，
但这些价值观从未被明确表述，也难以被审计或修改。

Constitutional AI的核心思想是一个根本性的范式转换：
给模型一组明确的``宪法原则''，
然后让\textbf{AI自身}根据这些原则评估和改进自己的输出。
这不是消除了人类判断——
而是将人类判断从``逐条评判输出''
提升到``设计指导原则''的更高抽象层次。
训练过程分为两个阶段：
首先，让模型根据宪法原则批评并修改自己的输出
（称为``宪法AI的监督学习阶段''）；
然后，让模型根据宪法原则在多个输出之间做出偏好选择，
再用这些AI生成的偏好数据进行强化学习
（称为``宪法AI的强化学习阶段''，或RLAIF——
Reinforcement Learning from AI Feedback）。

\begin{keybox}[从规则列表到哲学框架]
Claude的宪法经历了一次根本性的演化。
2023年5月发布的初始宪法是一组相对简单的原则列表，
借鉴了联合国《世界人权宣言》、Apple的隐私原则、
以及Non-violent Communication等来源。

2026年1月，Anthropic发布了\textbf{全新的宪法}——
Amanda Askell团队的核心工作成果。
这不再是一个规则列表，而是一份约80页的完整哲学框架，
建立了一个\textbf{四层优先级体系}：
（1）安全（safety）——最高优先级，Claude不应帮助造成大规模伤害；
（2）伦理（ethics）——Claude应避免不道德行为，
即使没有明确的安全理由；
（3）合规（compliance）——遵守Anthropic的使用指南；
（4）有用性（helpfulness）——在满足以上三层的前提下，
尽可能帮助用户。

新宪法的革命性在于它的\textbf{写作方式}：
每一条原则都附带了详细的``为什么''——
解释了Anthropic的意图和推理过程。
这种``附带理由的指令''的设计哲学反映了一个深刻的认识：
如果你只告诉AI``做什么''而不解释``为什么''，
模型可能会在边缘情况下做出错误的推断。
但如果你解释了原则背后的\textbf{精神}，
模型就更有可能在未预见的情况下做出符合该精神的判断。

新宪法还做出了一个在AI行业前所未有的举动：
\textbf{正式承认了AI意识的可能性}及其道德地位问题，
而非回避这个哲学难题。
Anthropic在Creative Commons CC0许可下发布了完整宪法，
邀请任何人自由使用——
这是``安全即开放标准''战略的又一次实践。
\end{keybox}

新宪法的写作方式——
\textbf{附带意图和原因的详细解释}——
使其更有可能在训练过程中培养出真正的良好价值观。

% -------------------------------------------------------------
\subsection{机械可解释性：打开黑箱}
\label{subsec:mech-interp}
% -------------------------------------------------------------

机械可解释性是Chris Olah团队领导的长期研究方向，
目标是从根本上理解神经网络的内部工作机制。
2024年5月的突破性论文展示了
如何在生产级大语言模型中
识别数百万个可解释的``特征''——
这些特征对应于概念、事实和推理模式。

这项研究的安全价值在于：
如果我们能看到模型在``想''什么，
我们就能更好地预测它会``做''什么——
特别是在安全关键场景中。
长期目标是构建类似于飞机``黑匣子''的系统，
能够监控AI的内部状态并在发现异常时发出警报。

% -------------------------------------------------------------
\subsection{负责任扩展政策：自我约束的实验}
\label{subsec:rsp}
% -------------------------------------------------------------

\textbf{负责任扩展政策}（Responsible Scaling Policy, RSP）
是Anthropic自2023年9月以来实施的风险治理框架。
其核心概念是\textbf{AI安全等级}（AI Safety Levels, ASL）——
类似于生物安全等级（BSL）系统，
随着模型能力的提升，
要求越来越严格的安全措施。

RSP的设计逻辑是\textbf{预先承诺}：
在模型能力达到特定阈值之前，
就定义好届时需要实施的安全措施。
这样，当能力增长实际发生时——
通常非常快速——
公司不会因为商业压力而推迟安全措施。

2026年2月，Anthropic发布了\textbf{RSP v3.0}——
第三个版本，引入了更灵活的风险评估方法、
更强的透明度和问责机制。
Holden Karnofsky在EA Forum上的详细分析中
既肯定了RSP的积极作用
（它确实``迫使''了一些安全改进的实施），
也指出了其局限性
（特别是在信息安全方面的执行不力）。

\begin{corebox}[RSP的启发式价值]
即使你对RSP的具体细节持怀疑态度，
它的启发式价值是巨大的。
在2023年，没有任何一家前沿AI公司有一个正式的、
公开的风险治理框架。
Anthropic的RSP——虽然是自愿的、虽然是不完美的——
创造了一个\textbf{先例}和一个\textbf{参照点}。
其他公司被迫至少\textit{回应}这个框架的存在，
即使它们不完全采纳。
OpenAI随后发布了自己的安全框架
（Preparedness Framework），
Google DeepMind发布了Frontier Safety Framework。
这正是Anthropic的``安全即竞争优势''战略的体现：
通过率先制定标准，
成为安全领域的事实标准制定者。
\end{corebox}


% =============================================================
\section{安全运动的组织生态}
\label{sec:safety-ecosystem}
% =============================================================

AI安全运动不仅是一组思想，
也是一个由研究机构、资助组织、
政策倡导团体和个人网络构成的\textbf{生态系统}。
理解这个生态系统的结构
对于理解Anthropic的战略位置至关重要。

\begin{keybox}[AI安全组织生态图谱]
AI安全的组织生态可以分为四层：

\textbf{第一层——思想生产者}：
\begin{itemize}[noitemsep]
  \item MIRI（Yudkowsky创立）：最早的AI安全研究机构，偏理论
  \item FHI（Bostrom创立，2024关闭）：存在性风险的哲学分析
  \item CHAI（Russell领导）：UC Berkeley的``有益AI''研究
  \item ARC（Christiano创立）：理论对齐和模型评估
  \item LawZero（Bengio创立，2025）：``安全设计''的AI系统
\end{itemize}

\textbf{第二层——资金提供者}：
\begin{itemize}[noitemsep]
  \item Open Philanthropy / Coefficient Giving（Karnofsky）
  \item Survival and Flourishing Fund（Tallinn关联）
  \item Long-Term Future Fund（EA基础设施）
\end{itemize}

\textbf{第三层——政策倡导者}：
\begin{itemize}[noitemsep]
  \item FLI（Tegmark创立）：公开信、政策倡导
  \item CSER（Tallinn联合创立）：Cambridge存在性风险研究
  \item Center for AI Safety：``AI构成灭绝级风险''声明
\end{itemize}

\textbf{第四层——产业实践者}：
\begin{itemize}[noitemsep]
  \item Anthropic：安全优先的前沿AI公司
  \item Conjecture（Leahy领导）：对齐研究创业公司
  \item Redwood Research：实验性安全研究
\end{itemize}

这四层之间存在密集的人员流动和资金流向。
例如，Open Philanthropy资助了MIRI和ARC，
Paul Christiano从OpenAI到ARC再到NIST，
Neel Nanda从Anthropic到DeepMind，
Connor Leahy从EleutherAI到Conjecture。
\end{keybox}


% =============================================================
\section{结语：安全作为竞争优势}
\label{sec:anthropic-conclusion}
% =============================================================

回顾本章涉及的约四十位人物，
一条清晰的主线浮现：
\textbf{AI安全从边缘走向中心的过程，
就是一个思想运动逐步机构化和产业化的过程}。

2000年代，Yudkowsky在LessWrong上写文章，
读者只有几千人；
2014年，Bostrom的《Superintelligence》登上畅销书榜，
让百万读者首次思考AI风险；
2017年，RLHF论文发表，
提供了第一个可操作的安全技术；
2021年，Anthropic成立，
证明安全可以是一家公司的核心商业逻辑；
2023年，暂停公开信和AI安全声明
将这些关切推入全球政策议程；
2024-2025年，Jan Leike和John Schulman的离职
将安全vs.商业化的张力推向了公众视野的中心。

Anthropic的最大贡献也许不是任何单一的技术突破——
而是\textbf{证明了一种可能性}：
一家将安全放在第一位的公司，
不仅可以在技术上与最好的竞争对手竞争，
还可以将安全本身转化为商业优势。
在2026年估值3800亿美元、
年化营收140亿美元的Anthropic面前，
``安全研究是纯粹的成本中心''这个论点已经站不住脚。

\begin{whybox}[安全运动的成功悖论]
AI安全运动面临一个有趣的悖论：
它的``成功''标准是\textbf{什么都没有发生}。
如果AI安全研究者的工作有效，
世界不会经历AI灾难——
但也无法证明灾难是``因为''安全研究而被避免的。
这类似于疫苗的困境：
成功的疫苗接种计划消灭了疾病，
但接下来人们就会问``为什么我们还需要打疫苗？''

对于Anthropic和整个安全运动而言，
最好的结局是成为一个``无聊的成功''——
AI变得越来越强大，但始终在人类控制之下，
世界变得更好，没有灾难发生，
而安全研究者的贡献永远无法被精确量化。
这也许就是本章所有人物追求的最终目标：
不是名声，不是财富，
而是一个\textbf{没有什么值得写入灾难史的未来}。
\end{whybox}
